\documentclass[12pt]{article}
\usepackage[utf8]{inputenc}
\usepackage[T1]{fontenc}
\usepackage{hyperref}
\title{Todrah Language Resources for K-12 Curriculum Development}
\author{Compiled by Claude Sonnet 5 (claude-sonnet-5) for Andr\'as Kornai}
\date{}
\begin{document}
\maketitle
\section{Introduction}
Todrah (ISO 639-3 tdr) is an Austroasiatic (North Bahnaric) language of Kon Tum Province, Vietnam, with about 6,000 speakers (Wikipedia, 2007) or about 6,900 (Joshua Project). Its two dialects, Sodrah and Xodrah, are quite distinct. Speakers are officially classified as Sedang. Alternate names include Sedang Didrah and Xo-dang.
\section{Primary Educational Resources}
No primers, textbooks or school programmes for Todrah were found.
\section{Linguistic and Reference Materials}
\begin{itemize}
\item Wikipedia, Todrah language: \url{https://en.wikipedia.org/wiki/Todrah_language}
\item Gregerson and Smith 1973, on the development of the Todrah register (Mon-Khmer Studies 4), bibliographic record: \url{https://glottolog.org/resource/reference/id/157626}
\end{itemize}
\section{Religious and Community Resources}
\begin{itemize}
\item Joshua Project, Todrah in Vietnam: \url{https://joshuaproject.net/people_groups/15526}
\end{itemize}
\section{Digital and Audio Resources}
None found.
\section{Curriculum Development Context}
Little is documented. Because speakers are administratively counted as Sedang, any materials may need to be coordinated with Sedang-language education in Vietnam, which was not examined in this search.
\end{document}